\documentclass{article}
\usepackage[utf8]{inputenc}
\usepackage[T1]{fontenc}
\usepackage{geometry}
\geometry{
    a4paper,
    left=2.5cm,
    right=2.5cm,
    top=2.5cm,
    bottom=2.5cm
}
\usepackage{graphicx}
\usepackage{float}
\usepackage{hyperref}
\usepackage[english]{babel}
\usepackage{amsmath}
\usepackage{amssymb}
\usepackage{array}
\usepackage{booktabs}
\usepackage{caption}
\usepackage{subcaption}
\usepackage{enumitem}
\usepackage{fancyhdr}
\usepackage{titlesec}
\usepackage{url}
\usepackage{textcomp}
\usepackage{xcolor}
\usepackage{listings}
\usepackage{multirow}
\usepackage{longtable}
\usepackage{tabularx}
\usepackage{times}

% Set up page style
\pagestyle{fancy}
\fancyhf{}
\rfoot{\thepage}
\renewcommand{\headrulewidth}{0pt}

% Title formatting
\titleformat{\chapter}[display]{\normalfont\bfseries}{\chaptertitle\thechapter}{0pt}{\Huge}
\titlespacing*{\chapter}{0pt}{50pt}{40pt}

% Document info
\title{Template Document}
\author{Academic Support Agent}
\date{\today}

\begin{document}
\maketitle

\maketitle

\begin{resumo} 
O artigo devem possuir obrigatoriamente um resumo, no seu idioma original, com 100 a 200 palavras ou 10 linhas no m ximo; seguido de um conjunto de palavras-chaves (de 3 at  5, escritas em min sculas), tamb m no idioma original do artigo. Logo a seguir deve aparecer a tradu o do resumo e das palavras-chaves, segundo o seguinte crit rio: se o idioma do artigo for o portugu s ou espanhol, a tradu o ser  para o ingl s; se o idioma do artigo for o ingl s, a tradu o ser  para o portugu s. (Abstract, Keywords); se for espanhol (Resumen) ou ingl s (Abstract), a tradu o ser  para o portugu s. O resumo deve ser claro, sucinto e explicar, em um  nico par grafo, o(s) objetivo(s) pretendido(s), procurando justificar sua import ncia (sem incluir refer ncias bibliogr ficas), os principais procedimentos adotados, os resultados mais expressivos e conclus es.
\end{resumo}

\begin{palavas-chave}
    Palavra Chave 1, Palavra Chave 2, Palavra Chave 3, Palavra Chave 4, Palavra Chave 5.
\end{palavas-chave}

\begin{englishtitle}
    English Title (Renote Journal)
\end{englishtitle}

\begin{abstract}
O artigo devem possuir obrigatoriamente um resumo, no seu idioma original, com 100 a 200 palavras ou 10 linhas no m ximo; seguido de um conjunto de palavras-chaves (de 3 at  5, escritas em min sculas), tamb m no idioma original do artigo. Logo a seguir deve aparecer a tradu o do resumo e das palavras-chaves, segundo o seguinte crit rio: se o idioma do artigo for o portugu s ou espanhol, a tradu o ser  para o ingl s; se o idioma do artigo for o ingl s, a tradu o ser  para o portugu s. (Abstract, Keywords); se for espanhol (Resumen) ou ingl s (Abstract), a tradu o ser  para o portugu s. O resumo deve ser claro, sucinto e explicar, em um  nico par grafo, o(s) objetivo(s) pretendido(s), procurando justificar sua import ncia (sem incluir refer ncias bibliogr ficas), os principais procedimentos adotados, os resultados mais expressivos e conclus es.
\end{abstract}

\begin{keywords}
    Palavra Chave 1, Palavra Chave 2, Palavra Chave 3, Palavra Chave 4, Palavra Chave 5.
\end{keywords}

\section{Estrutura}

Tamanho da folha: para  impress o em papel  A4 (21,0 cm x 29,7 cm) sendo observadas as seguintes margens: superior 3,0; cm inferior 2,0 cm; esquerda e direita 3,0 cm. Alinhar a margem direita evitando separa es sil bicas com barras ou outros sinais.

Utilizar fonte "Times" tamanho 12 ou equivalente. N o utilizar fontes que ocupam muito espa o tal como "Bookman".  No texto utilizar espa o 1.  O fim de uma se o e o cabe alho da pr xima s o separados por um espa o extra.

Todas as p ginas do trabalho devem ser numeradas, com exce o da folha de rosto.  Os n meros, em algarismos ar bicos, s o colocados, no centro da margem superior. Sugere-se utilizar o comando "cabe alho" do editor de texto para especificar a pagina o. Entre a pagina o e o texto deixar algum espa o para destacar os dois elementos. 

\section{Teste}
O n mero de tabelas e figuras utilizadas no artigo deve ser limitado a compreens o e elucida o do texto. Devem ser inseridas no corpo do texto, para identifica o da sua posi o e do tamanho aproximado (Figura~\ref{fig:exampleFig1}).

\subsection{Subse o}
O n mero de tabelas e figuras utilizadas no artigo deve ser limitado a compreens o e elucida o do texto. Devem ser inseridas no corpo do texto, para identifica o da sua posi o e do tamanho aproximado (Figura~\ref{fig:exampleFig1}).

\section{Figures, Tables and Captions}\label{sec:figs}

O n mero de tabelas e figuras utilizadas no artigo deve ser limitado a compreens o e elucida o do texto. Devem ser inseridas no corpo do texto, para identifica o da sua posi o e do tamanho aproximado (Figura~\ref{fig:exampleFig1}).

Tabelas e figuras possuem numera o independente, que deve ser feita seq encialmente na ordem em que s o citadas no texto. Devem tamb m ter uma legenda auto-explicativa, sendo que as tabelas ter o legendas na parte superior e as figuras as ter o na parte inferior, centralizadas em rela o   tabela ou figura. A legenda inicia com o termo "Tabela" ou "Figura" (primeira letra em mai scula), de acordo com o caso, seguido de um espa o e do n mero de ordem seq encial, em algarismos ar bicos, seguido de h fen entre espa os e do texto da legenda, com a primeira letra da primeira palavra em mai scula e as demais em min scula, exce o feitas  quelas que normalmente s o escritas em mai sculas. Devem ser citadas no texto como "Tabela \ref{tab:exTable1}" e "Figura" seguidas de espa o e do n mero correspondente.

\begin{figure}[ht]
\centering
\includegraphics[width=.5\textwidth]{fig1.jpg}
\caption{A typical figure}
\label{fig:exampleFig1}
\end{figure}

\begin{figure}[ht]
\centering
\includegraphics[width=.3\textwidth]{fig2.jpg}
\caption{This figure is an example of a figure caption taking more than one
  line and justified considering margins mentioned in Section~\ref{sec:figs}.}
\label{fig:exampleFig2}
\end{figure}

\begin{table}[ht]
\centering
\caption{Variables to be considered on the evaluation of interaction
  techniques}
\label{tab:exTable1}
\includegraphics[width=.7\textwidth]{table.jpg}
\end{table}

\section{Agradecimentos}
Usados apenas quando necess rio. Devem ser sucintos, diretos e dirigidos apenas a pessoas ou institui es que contribu ram substancialmente para a elabora o do trabalho.

\section{Refer ncias}

As refer ncias a autores e fontes s o inseridas no texto colocando entre par nteses o sobrenome do autor, com inicial mai scula, seguida da data, conforme o exemplo: \cite{smith:99}. Quando for conhecida a pagina o, pode-se inclu -la, por exemplo, (Tarouco, 1991, p. 237). Havendo mais de um t tulo do(s) autor(es) no mesmo ano, deve-se distingui-las utilizando uma letra min scula (a,b,c) depois da data, por exemplo, (Tarouco, 1991b). Quando houver tr s ou mais autores, no texto ser  citado apenas o primeiro autor seguido de "et al.", mas nas refer ncias bibliogr ficas, no final do artigo, os demais nomes tamb m dever o aparecer. Quando o nome do autor   citado diretamente no texto pode-se colocar entre par nteses apenas a data de publica o, por exemplo,

\citeonline{boulic:91}. Na cita o de cita o, identifica-se a obra diretamente consultada; o autor e/ou a obra citada nesta   assim indicado: \cite{knuth:84} citada por Vicari \cite{teste}.

As refer ncias completas de cada autor e fonte citadas no texto devem aparecer no final do artigo sob o t tulo "Refer ncias Bibliogr ficas", ordenadas alfabeticamente pelo(s) sobrenomes(s) do(s) autore(s).

Os autores do artigo devem certificar-se de que as refer ncias citadas no texto contam da lista de refer ncias com datas exatas e nomes de autores corretamente grafados. A exatid o dessas refer ncias   de responsabilidade dos autores do artigo. Comunica es pessoais, trabalhos in ditos ou em andamento poder o ser citados quando absolutamente necess rios, mas n o devem ser inclu dos na lista de refer ncias bibliogr ficas; apenas citados no texto ou nas Notas do texto.

\addto{\captionsbrazil}{% caso use \usepackage[brazil]{babel}
\renewcommand{\bibname}{Refer\^{e}ncias Teses}
}
\end{document}